\documentclass{report}
\usepackage{amsmath,amssymb}
\setlength{\parindent}{0mm}
\setlength{\parskip}{1em}
\begin{document}
\begin{center}
\rule{6in}{1pt} \
{\large Jacob B. Schroder \\
{\bf Multigrid Reduction in Time: Recent theoretical results}}

Lawrence Livermore National Laboratory
\\
{\tt schroder2@llnl.gov}\\
R. Falgout\\
Tz. Kolev\\
U. Yang\\
V. Dobrev\\
A. Petersson\end{center}

The need for parallel-in-time is being driven by changes in computer
architectures, where future speedups will be available through greater
concurrency, but not faster clock speeds, which are stagnant. This leads
to a bottleneck for sequential time marching schemes, because they lack
parallelism in the time dimension. In this talk, we examine an
optimal-scaling parallel time integration method,
multigrid-reduction-in-time (MGRIT). MGRIT applies multigrid to the time
dimension by solving the (non)linear systems that arise when solving for
multiple time steps simultaneously. The result is a versatile approach
that is nonintrusive and wraps existing time evolution codes. In this
talk, we present the MGRIT framework and then discuss some recent
theoretical convergence results. Some typical simplifying assumptions are
made, such as a two-grid method, uniform time-line and constant
coefficient PDE. The convergence estimates are then explored in a variety
of settings, including common parabolic and hyperbolic spatial
discretizations coupled with implicit Runge-Kutta time-stepping schemes.
Overall, the convergence estimates are sharp when compared to numerical
experiments and are good predictors of multilevel results.


\end{document}
